\documentclass[12pt,a4paper,twoside,openright]{book}
%\usepackage[work]{MCthesis}            % you can select this to work

\usepackage[final]{MCthesis}            % this is for the final submission

%%%%%%%%%%%%%%%%%%%%%%%%%%%%%%%%%%%%%%%%%%%%%%%%%%%%%%
%%%%%%%%%%%%%%%%%%%%%%%%%%%%%%%%%%%%%%%%%%%%%%%%%%%%%%
\usepackage[utf8]{inputenc}         % allow utf-8 input
\usepackage[T1]{fontenc}            % use 8-bit T1 fonts
\usepackage{hyperref}               % hyperlinks

% Hiding the red boxes around links
\hypersetup{
    hidelinks
}

\usepackage{url}                    % simple URL typesetting
\usepackage{booktabs}               % professional-quality tables
\usepackage{nicefrac}               % compact symbols for 1/2, etc.
\usepackage{microtype}              % microtypography
\usepackage{xcolor}                 % colors
\usepackage{amssymb,amsmath,amsfonts,amsbsy, amsthm}
\usepackage{mathrsfs}
\usepackage{bm}                     % bold vectors
\usepackage{listings}               % to include R-Code
\usepackage[section]{placeins}      % forces floats to appear in the section they are placed in
\usepackage[]{apacite}              % bibliography example
\usepackage[toc,page]{appendix}
\usepackage{framed}
\usepackage{float}
\usepackage{algorithm}              % For your pseudo-code
\usepackage{algpseudocode}          % For your pseudo-code
%%%%%%%%%%%%%%%%%%%%%%%%%%%%%%%%%%%%%%%%%%%%%%%%%%%%%%
%%%%%%%%%%%%%%%%%%%%%%%%%%%%%%%%%%%%%%%%%%%%%%%%%%%%%%
\usepackage{blindtext}

%%%%%%%%%%%%%%%%%%%%%%%%%%%%%%%%%%%%%%%%%%%%%%%%%%%%
% SET VARIABLES FOR TITLE PAGE
\title{Optimal Transport in linear Independent Component Analysis}
\author{Ashutosh Jha}
\fromCity{Tuebingen}
\matnr{6639615}
\date{\today} % submission date

%%%%%%%%%%%%%%%%%%%%%%%%%%%%%%%%%%%%%%%%%%%%%%%%%%%%
%%%%%%%%%%%%%%%%%%%%%%%%%%%%%%%%%%%%%%%%%%%%%%%%%%%%
% SET VARIABLES FOR SUPERVISION AND EXAMINERS
\supervisor{Prof. Dr. Joachim Grammig\textsuperscript{1}}
           {Prof. Dr. Michel Besserve\textsuperscript{2,3}}
           {Dr. Simon Buchholz\textsuperscript{2}}

\reviewer{Prof. Dr. Joachim Grammig\textsuperscript{1}}
         {Prof. Dr. Augustin Kelava\textsuperscript{4}
          \\[1.5cm] % Adds space before the affiliations
          \multicolumn{2}{l}{\small \textsuperscript{1} Chair of Statistics, Econometrics and Empirical Economics, University of Tübingen} \\
          \multicolumn{2}{l}{\small \textsuperscript{2} Empirical Inference, Max Planck Institute for Intelligent Systems, Tübingen} \\
          \multicolumn{2}{l}{\small \textsuperscript{3} Institute for Machine Learning and AI, TU Braunschweig} \\
          \multicolumn{2}{l}{\small \textsuperscript{4} Methods Center, University of Tübingen}
         }
         {}
%%%%%%%%%%%%%%%%%%%%%%%%%%%%%%%%%%%%%%%%%%%%%%%%%%%%

%%%%%%%%%%%%%%%%%%%%%%%%%%%%%%%%%%%%%%%%%%%%%%%%%%%%%%%%%%%%%%%%%%%

\begin{document}

% The \maketitle command in your MCthesis package automatically handles 
% the title page, declaration, \frontmatter, \tableofcontents, and \mainmatter.
\maketitle

% -----------------------------------------------------------------
% CHAPTER 1: INTRODUCTION
% -----------------------------------------------------------------
\chapter{Introduction}

\section{Motivation: Disentangling Latent Variables}
Blind Source Separation, specifically formulated as Independent Component Analysis (ICA), is a fundamental problem in modern statistics \cite{jutten1985}. 
It has widespread applications across diverse fields, including machine learning, signal processing, econometrics, and neuroscience. 
The primary objective of ICA is to recover hidden, statistically independent sources from an observed set of mixtures. 
In this thesis, we focus on the linear mixture setting. ICA is frequently contrasted with Principal Component Analysis (PCA); 
however, while PCA merely ensures that the resulting components are uncorrelated, ICA enforces the much stronger condition of full statistical independence \cite{hyvarinen2001}. 
This rigorous criterion allows ICA to successfully disentangle complex latent variables that PCA cannot fully separate.

\section{Intuition from Central Limit Theorem (CLT)}
For the linear ICA model to be identifiable, a strict condition must be met: at most one of the underlying independent sources can follow a Gaussian distribution. 
If multiple sources are Gaussian, any linear combination of them results in another Gaussian distribution, 
rendering the original signals rotationally symmetric and impossible to uniquely separate. 
This fundamental limitation naturally translates the challenge of finding independent components into an optimization problem: we must find the projections that maximize non-Gaussianity.

The theoretical justification for this approach is deeply rooted in the Central Limit Theorem (CLT). 
The CLT establishes that the sum or any linear combination of independent random variables tends toward a Gaussian distribution, 
regardless of the distributions of the original variables. Consequently, any observed mixture of independent, 
non-Gaussian sources will inherently appear more Gaussian than the constituent sources themselves. 
Therefore, by extracting the components that are the least Gaussian, we are effectively isolating the original independent signals.

\section{Information Theory, Thermodynamics and CLT}
The use of non-gaussianity as a measure for independence of components is also motivated by information 
theory. For a continuous random variable with a fixed variance, it is a foundational result that the Gaussian 
distribution uniquely possesses the maximum possible Shannon entropy. This mathematical property reveals a 
profound intrinsic link between information theory, the Central Limit Theorem, and statistical mechanics. 

In thermodynamics, isolated physical systems naturally evolve toward states of maximum entropy, or maximum 
disorder. Viewed through this interdisciplinary lens, the Central Limit Theorem serves as the statistical 
manifestation of this physical principle. The mathematical process of linearly mixing independent sources 
inherently increases the overall entropy of the system, inevitably driving the joint distribution toward 
the high-entropy Gaussian state. 

Recovering the underlying independent components, then, is conceptually equivalent to reversing this 
thermodynamic process. It requires projecting the empirical data away from the "disordered" high-entropy 
Gaussian manifold, intentionally guiding it back toward the highly structured, lower-entropy state of the 
original independent signals.

\section{Problem Statement and Research Objectives}
In practice, the ICA pipeline begins by centering and whitening the observed mixture. 
Whitening is a crucial preprocessing step that removes linear correlations and scales the variances to unity. Once the data is whitened, 
the search for independent components is geometrically restricted to finding an optimal orthogonal rotation. 
Specifically, the problem reduces to searching the unit sphere of the Stiefel manifold, space of orthogonal matrices, 
for the rotation matrix that maximizes the non-Gaussianity of the projected data.

The primary novelty of this research lies in proposing and investigating Optimal Transport metric, the Wasserstein distance, as the objective function for this rotational search. 
Our objective is to find the orthogonal candidates on the unit sphere that maximize the Wasserstein distance to a standard normal distribution. 
While traditional ICA algorithms often rely on sample approximations of population moments (such as kurtosis) or surrogate information-theoretic measures (such as negentropy approximations), 
the Wasserstein metric offers a distinct advantage: 
it evaluates the true global geometry of the empirical distribution using order statistics (sorted data), providing a highly robust, non-parametric measure of distance to Gaussianity. 

\section{Structure of the Thesis}
The remainder of this thesis is structured as follows. Chapter 2 establishes the theoretical foundations of linear ICA, information geometry, and Optimal Transport. 
Chapter 3 formalizes the squared Wasserstein-ICA (W2-ICA) framework, proving that the squared Wasserstein (W2) distance rigorously bounds mixture non-Gaussianity. 
Chapter 4 details the algorithmic enhancements required to optimize this metric on the Stiefel manifold efficiently. 
Chapter 5 presents our experimental evaluation, demonstrating the failure modes of standard proxy optimization (such as FastICA) and proposing a robust ground cost modification to handle outliers. 
Finally, Chapter 6 concludes the thesis with a summary of findings and directions for future research.


% -----------------------------------------------------------------
% CHAPTER 2: THEORETICAL FOUNDATIONS
% -----------------------------------------------------------------
\chapter{Theoretical Foundations}

\section{Linear Independent Component Analysis (ICA)}

\subsection{Notation and the Linear Mixing Model}
The foundational assumption of the linear Independent Component Analysis model is that the observed signals are generated by a linear combination of statistically independent latent sources. 
Let the observed data be denoted by a random vector $\mathbf{X} \in \mathbb{R}^d$, and the latent independent sources by $\mathbf{S} \in \mathbb{R}^d$. 
The linear mixing model is formally defined as
\begin{equation}
    \mathbf{X} = \mathbf{A} \mathbf{S}
\end{equation}
where $\mathbf{A} \in \mathbb{R}^{d \times d}$ is an unknown, invertible mixing matrix \cite{hyvarinen2001}. Expanded into vector notation, the model represents the observed variables as
\begin{equation}
    \begin{bmatrix}
    x_1 \\
    x_2 \\
    \vdots \\
    x_d
    \end{bmatrix}
    = 
    \begin{bmatrix}
    a_{11} & a_{12} & \cdots & a_{1d} \\
    a_{21} & a_{22} & \cdots & a_{2d} \\
    \vdots & \vdots & \ddots & \vdots \\
    a_{d1} & a_{d2} & \cdots & a_{dd}
    \end{bmatrix}
    \begin{bmatrix}
    s_1 \\
    s_2 \\
    \vdots \\
    s_d
    \end{bmatrix}.
\end{equation}
The core objective of ICA is to estimate an unmixing matrix $\mathbf{B} \approx \mathbf{A}^{-1}$ such that the recovered components $\mathbf{Z} = \mathbf{B}\mathbf{X}$ are as 
statistically independent as possible, thereby reconstructing the original sources $\mathbf{S}$.

\subsection{Identifiability and Ambiguities}
For the linear ICA model to be identifiable, it is a strict mathematical requirement that at most one of the underlying sources $s_i$ follows a Gaussian distribution. 
If two or more sources are Gaussian, any orthogonal transformation of those sources yields another set of independent Gaussian variables. Consequently, 
the mixing matrix cannot be uniquely determined from the observed data. Furthermore, even when the non Gaussianity assumption is met, the ICA model contains two inherent ambiguities. 
First, the variances (energies) of the independent components cannot be determined, which leads to a scaling ambiguity. Second, the order of the sources is arbitrary, 
leading to a permutation ambiguity. These ambiguities are generally acceptable in practice because the structural shape and independence of the source distributions are preserved.

\subsection{Centering and Whitening}
To simplify the optimization problem, the observed data $\mathbf{X}$ is subjected to a two step preprocessing phase consisting of centering and whitening. 
Centering ensures that the data has a zero mean by subtracting the expectation
\begin{equation}
    \mathbb{E}[\mathbf{X}] = \mathbf{0}.
\end{equation}
Given centered data, the covariance matrix simplifies to $\mathrm{Cov}(\mathbf{X}) = \mathbb{E}[\mathbf{X}\mathbf{X}^\top]$. 
Whitening is then applied to transform the observed vector into a new vector $\widetilde{\mathbf{X}}$ whose components are uncorrelated and possess unit variance. 
This is achieved using the eigenvalue decomposition of the covariance matrix
\begin{equation}
    \mathrm{Cov}(\mathbf{X}) = \mathbf{V}\mathbf{\Lambda}\mathbf{V}^\top
\end{equation}
where $\mathbf{V}$ is the orthogonal matrix of eigenvectors and $\mathbf{\Lambda}$ is the diagonal matrix of eigenvalues $\lambda_j$. 
The whitening transformation matrix $\mathbf{W}$ is constructed by scaling the eigenvectors by the inverse square root of the eigenvalues
\begin{equation}
    \mathbf{W} = \mathbf{\Lambda}^{-1/2} \mathbf{V}^\top.
\end{equation}
Applying this transformation yields the whitened data $\widetilde{\mathbf{X}} = \mathbf{W}\mathbf{X}$. The covariance of this new variable is the identity matrix
\begin{equation}
    \mathrm{Cov}(\widetilde{\mathbf{X}}) = \mathbf{W} \mathbb{E}[\mathbf{X}\mathbf{X}^\top] \mathbf{W}^\top = \mathbf{\Lambda}^{-1/2} \mathbf{V}^\top \mathbf{V} \mathbf{\Lambda} \mathbf{V}^\top \mathbf{V} \mathbf{\Lambda}^{-1/2} = \mathbf{I}_d.
\end{equation}
After whitening, the search for the unmixing matrix is drastically simplified. 
If we define the unmixing operation on the whitened data as $\widetilde{\mathbf{Z}} = \mathbf{U} \widetilde{\mathbf{X}}$, the covariance of the recovered components is
\begin{equation}
    \mathrm{Cov}(\widetilde{\mathbf{Z}}) = \mathbf{U} \mathrm{Cov}(\widetilde{\mathbf{X}}) \mathbf{U}^\top = \mathbf{U} \mathbf{I}_d \mathbf{U}^\top = \mathbf{U}\mathbf{U}^\top.
\end{equation}
Because the independent components must have unit variance, we require $\mathbf{U}\mathbf{U}^\top = \mathbf{I}_d$. This reveals that the unmixing matrix $\mathbf{U}$ must be perfectly orthogonal. 
The ICA problem is thus reduced to finding an optimal rotation on the Stiefel manifold.

\subsection{The Insufficiency of Principal Component Analysis}
Principal Component Analysis is frequently used to decorrelate data, which corresponds to finding the matrix $\mathbf{V}^\top$ in the eigenvalue decomposition. 
However, making a sample decorrelated and whitened is not sufficient for source separation when the data is non Gaussian. 
Once the data is whitened such that $\mathbb{E}[\mathbf{X}\mathbf{X}^\top] = \mathbf{I}_d$, any rotational orientation satisfies the PCA criteria. 
For multivariate Gaussian distributions, decorrelation directly implies independence, leaving nothing further for PCA to accomplish. For non Gaussian data, however, 
higher order dependencies remain. ICA must therefore traverse the space of orthogonal matrices to find local optima of contrast functions such as $f(\mathbf{u}) = \mathbb{E}[(\mathbf{u}^\top \mathbf{X})^4]$ on the unit sphere $S_n$. 
Gradient descent and other optimization techniques are required to pinpoint the exact rotation that maximizes non Gaussianity and achieves full statistical independence.

\section{An Information Geometry Perspective on ICA}

\subsection{Product and Gaussian Manifolds}
Information geometry provides a rigorous framework for quantifying the relationship between independence and non Gaussianity by analyzing probability distributions as points on geometric manifolds. 
Following the perspective introduced by Cardoso \citeyear{cardoso2022}, ICA can be understood as projecting the empirical distribution onto two distinct subspaces. The first is the Product manifold $\mathcal{P}$, 
which contains all distributions where the components are strictly independent. The second is the Gaussian manifold $\mathcal{G}$, which encompasses all multivariate Gaussian distributions with zero mean and a fixed covariance structure. 
Standard decorrelation merely ensures that the covariance is diagonal, but full ICA operates at the intersection of these manifolds to find a linear transform that projects the empirical data as closely as possible onto the Product manifold \cite{cardoso2022}.

\subsection{Kullback-Leibler Pythagorean Identities}
Distances between these distributions are measured using the Kullback-Leibler (KL) divergence. For any random vector $Y$ with a joint density $P_Y$, 
the mutual information $\mathscr{I}(Y)$ measures the distance to the closest independent product distribution
\begin{equation}
    \mathscr{I}(Y) = D_{\mathrm{KL}}\left(P_Y \,\Vert \, \prod_{i} P_{Y_i}\right).
\end{equation}
A Pythagorean identity exists on the Product manifold. For any candidate product distribution $Q = \prod_{i} q_i$, the divergence splits precisely into the mutual information and the marginal divergences
\begin{equation}
    D_{\mathrm{KL}}\left(P_Y \Vert Q\right) = \mathscr{I}(Y) + \sum_{i} D_{\mathrm{KL}}\left(P_{Y_i} \Vert q_i\right).
\end{equation}
A similar Pythagorean identity holds for the Gaussian manifold $\mathcal{G}$. The KL divergence from $P_Y$ to any arbitrary Gaussian distribution $\mathcal{N}(\Sigma)$ decomposes into 
the divergence to the best Gaussian fit and the divergence between the two Gaussian models
\begin{equation}
    D_{\mathrm{KL}}\big(P_Y \Vert \mathcal{N}(\Sigma)\big) = D_{\mathrm{KL}}\big(P_Y \Vert \mathcal{N}(\operatorname{Cov} Y)\big) + D_{\mathrm{KL}}\big(\mathcal{N}(\operatorname{Cov} Y) \Vert \mathcal{N}(\Sigma)\big).
\end{equation}

\subsection{Relating Independence, Correlation, and Non-Gaussianity}
These geometric identities can be combined to relate independence directly to non Gaussianity. We define the non Gaussianity of the joint variable $Y$ as $G(Y) = D_{\mathrm{KL}}(P_Y \,\Vert\, \mathcal{N}(\operatorname{Cov} Y))$, and the scalar measure of correlation as $C(Y) = D_{\mathrm{KL}}(\mathcal{N}(\operatorname{Cov} Y) \,\Vert\, \mathcal{N}(\operatorname{diag\,Cov} Y))$. Cardoso's unified theorem states that
\begin{equation}
    \mathscr{I}(Y) + \sum_{i} G(Y_i) = C(Y) + G(Y).
\end{equation}
Crucially, the joint non Gaussianity $G(Y)$ is invariant under invertible linear transformations. Furthermore, after the data has been whitened, 
the covariance matrix is the identity matrix, meaning $\operatorname{Cov}(Y) = \operatorname{diag\,Cov}(Y)$ and therefore the correlation term $C(Y) = 0$. 
The relationship then drastically simplifies to
\begin{equation}
    \mathscr{I}(Y) = G(Y) - \sum_{i} G(Y_i).
\end{equation}
Because $G(Y)$ is constant under rotation, minimizing the mutual information $\mathscr{I}(Y)$ is mathematically equivalent to maximizing the sum of the 
marginal non Gaussianities $\sum_i G(Y_i)$. This rigorously proves that centering and whitening constrain the data to a subspace where true statistical independence can be 
achieved solely by maximizing non Gaussianity \cite{cardoso2022}.

\section{Optimal Transport Theory}

\subsection{The Monge-Kantorovich Problem}
While information geometry utilizes Kullback-Leibler divergence, Optimal Transport offers an alternative metric space that is highly sensitive to the underlying geometry of the data. 
First formulated by Gaspard Monge \citeyear{monge1781}, the problem was originally conceptualized as finding the most cost effective way to move a distribution of mass, such as a pile of soil, 
to a target destination. 

Mathematically, this is framed as finding a transport map $T$ that pushes forward a source probability measure $\mu$ to a target measure $\nu$, denoted as $T_{\#} \mu = \nu$. 
This means that for any measurable set $A$, the mass assigned by $\nu$ is equal to the mass assigned by $\mu$ to the preimage of $A$, such that $\nu(A) = \mu(T^{-1}(A))$. 
Because finding a deterministic map $T$ is highly restrictive, Kantorovich \citeyear{kantorovich1942} later relaxed this formulation by introducing transport plans. This allowed mass from a single source 
location to be split and distributed across multiple targets. This relaxation transformed the rigorous mapping requirement into a joint probability distribution problem, 
making the computation of transport costs mathematically tractable.

\subsection{The Earth Mover's Distance ($W_1$)}
The minimal cost required to transform one probability distribution into another is formalized as the Wasserstein distance \cite{villani2003}. The $W_1$ distance, commonly referred to as the Earth Mover's Distance, 
quantifies the absolute effort of this transformation. For one dimensional distributions with cumulative distribution functions $F$ and $G$, the $W_1$ distance is computed as the integral of the 
absolute differences between the cumulative functions
\begin{equation}
    W_1(\mu, \nu) = \int_{-\infty}^\infty |F(x) - G(x)| \, dx.
\end{equation}
Intuitively, $W_1$ sums the vertical area between the two distribution curves. While this metric is robust to extreme outliers, it often underestimates structural deviations in variance.

\subsection{The Squared Wasserstein Distance ($W_2$)}
To better capture variations in spread and central location, the Squared Wasserstein distance is utilized. It is most conveniently defined using the quantile functions, 
which are the inverses of the cumulative distribution functions $F^{-1}$ and $G^{-1}$. The $W_2$ distance is formulated as
\begin{equation}
    W_2(\mu, \nu) = \left( \int_0^1 |F^{-1}(t) - G^{-1}(t)|^2 \, dt \right)^{1/2}.
\end{equation}
By squaring the transport cost, the $W_2$ metric heavily penalizes local geometric distortions. When evaluating the distance between an empirical projection and a 
standard normal distribution $\mathcal{N}(0,1)$, $W_2$ provides a highly sensitive and reliable measure of Gaussianity. This makes it an ideal contrast function for Independent Component Analysis, 
as it maps the exact geometric deformations required to transform the candidate components into Gaussian noise.

\subsection{Brenier's Theorem and Caffarelli's Regularity}
The use of the squared Euclidean cost function in the $W_2$ distance provides profound geometric properties that are not present in $W_1$. 
Specifically, Brenier's Theorem establishes a direct link between the squared Wasserstein distance and convex analysis \cite{brenier1991}. 
The theorem states that if the source measure $\mu$ is absolutely continuous with respect to the Lebesgue measure, the optimal transport plan is uniquely given by a deterministic Monge map $T$. 
Crucially, Brenier proved that this optimal map is the gradient of a convex function $\phi$, such that $T = \nabla \phi$. 

The analytical behavior of this transport map is governed by Caffarelli's regularity theory \cite{caffarelli1992}. Caffarelli demonstrated that if both the source and target probability densities are 
smooth and bounded strictly away from zero on their supports, the convex potential $\phi$ is strictly convex and its gradient $T$ is continuously differentiable. 
This regularity implies that the optimal transport map does not tear or fold the probability space abruptly. 
In the context of linear ICA, the existence of a unique, smooth, and differentiable transport map mapping empirical projections to a Gaussian target provides the 
necessary mathematical foundation to analyze optimization gradients and theoretically bound mixture non-Gaussianities.

% -----------------------------------------------------------------
% CHAPTER 3: THE SQUARED WASSERSTEIN-ICA (W2-ICA) FRAMEWORK
% -----------------------------------------------------------------
\chapter{The Squared Wasserstein-ICA (W2-ICA) Framework}

\section{Formulating ICA as an Optimal Transport Problem}
As established in the preceding chapter, the process of centering and whitening restricts the search for independent 
components to the manifold of orthogonal matrices. Furthermore, Cardoso's unified theorem \cite{cardoso2022} demonstrates 
that within this whitened space, the mutual information of the projections is minimized precisely when the sum of their 
marginal non-Gaussianities is maximized. Traditional Independent Component Analysis algorithms approach this optimization 
by utilizing proxy contrast functions, such as kurtosis or negentropy approximations. While computationally inexpensive, 
these proxies often capture only specific statistical moments, ignoring the broader geometric structure of the underlying probability distributions.

This thesis proposes a paradigm shift by formulating ICA directly as an Optimal Transport problem. Instead of relying on 
proxy moments, we utilize the squared Wasserstein distance, $W_2^2$, as a direct, geometrically rigorous measure of non-Gaussianity. 
Let $\mathbf{X}$ be the whitened observed mixture, and let $\mathbf{u} \in S_n$ be a candidate projection vector on the unit sphere. 
The empirical distribution of the projected data is denoted by $\mathbb{P}_{\mathbf{u}^\top \mathbf{X}}$. We formally define the 
objective of the W2-ICA framework as finding the orthogonal rotation matrix whose column vectors $\mathbf{u}$ maximize the 
transport cost to a standard normal distribution $\mathcal{N}(0,1)$:
\begin{equation}
    \hat{\mathbf{u}} = \arg\max_{\mathbf{u} \in S_n} W_2^2\big(\mathbb{P}_{\mathbf{u}^\top \mathbf{X}}, \mathcal{N}(0,1)\big).
\end{equation}
By framing the problem in this manner, we leverage the Optimal Transport metric's unique ability to capture global geometric 
deviations from Gaussianity, providing a highly sensitive contrast function.

\section{Theoretical Motivation: Bounding Mixture Non-Gaussianity}
To mathematically justify the use of $W_2^2$ as an objective function for ICA, we must prove that maximizing this distance 
recovers the true independent sources. By the Central Limit Theorem, we possess the intuitive understanding that a mixture of 
non-Gaussian sources is more Gaussian than the constituent sources themselves. In the language of Optimal Transport, this implies 
that the Wasserstein distance between a mixture and a Gaussian distribution must be strictly bounded by the distances of its independent components. 


\subsection{Constructing the Common Source Coupling}
To formally prove this bound, we construct a specific joint probability distribution, or coupling, using a common source of randomness. 
Let $\mathbf{N} = (N_1, \dots, N_d)^\top$ be a vector of $d$ independent standard normal variables, such that $\mathbf{N} \sim \mathcal{N}(0, \mathbf{I}_d)$. 
For each true independent latent source $Z_i$, there exists an Optimal Transport map $T_i$ that pushes forward the standard normal distribution 
to the source distribution, denoted as $(T_i)_{\#} \mathcal{N}(0,1) = Z_i$. 

We can construct a random variable $X$ representing a linear mixture of the sources using a weight vector $\mathbf{w}$ defined on the 
unit sphere ($\sum_{i=1}^d w_i^2 = 1$). By applying the transport maps to our common noise vector, the mixture is defined as
\begin{equation}
    X = \sum_{i=1}^d w_i T_i(N_i).
\end{equation}
Because $T_i(N_i)$ follows the distribution of $Z_i$, it follows that $X$ perfectly models the weighted mixture $\mathbf{w} \cdot \mathbf{Z}$. 
To compare this mixture to a Gaussian distribution, we construct a target variable $Y$ using the exact same underlying noise vector $\mathbf{N}$ and the same weights:
\begin{equation}
    Y = \sum_{i=1}^d w_i N_i.
\end{equation}
Given that the weights sum to unity and the components $N_i$ are independent standard normals, the resulting variable $Y$ is guaranteed to follow 
a standard normal distribution, $Y \sim \mathcal{N}(0,1)$. The pair $(X, Y)$ therefore creates a valid coupling $\pi_{X,Y}^{\mathbf{N}}$ 
between the mixture distribution and the Gaussian target.

\subsection{Proof of the $W_2$ Upper Bound}
The squared Wasserstein distance is defined as the infimum of the expected squared cost over all possible valid couplings. 
Because the infimum represents the absolute optimal transport plan, the cost of the optimal plan must be less than or equal to the cost of our 
specifically constructed coupling $(X, Y)$. This yields the fundamental inequality:
\begin{equation}
    W_2^2(\mathbf{w} \cdot \mathbf{Z}, \mathcal{N}(0,1)) \le \mathbb{E}\left[|X - Y|^2\right].
\end{equation}
We evaluate the right side of this inequality by expanding the squared difference:
\begin{align}
    |X - Y|^2 &= \left| \sum_{i=1}^d w_i T_i(N_i) - \sum_{i=1}^d w_i N_i \right|^2 \nonumber \\
              &= \left( \sum_{i=1}^d w_i \big(T_i(N_i) - N_i\big) \right)^2.
\end{align}
Squaring the summation produces a combination of diagonal terms where $i=j$ and cross terms where $i \neq j$:
\begin{equation}
    |X - Y|^2 = \sum_{i=1}^d w_i^2 \big(T_i(N_i) - N_i\big)^2 + \sum_{i \neq j} w_i w_j \big(T_i(N_i) - N_i\big)\big(T_j(N_j) - N_j\big).
\end{equation}
Taking the mathematical expectation of this expansion drastically simplifies the expression. Because the underlying noise variables $N_i$ and $N_j$ are independent, 
and the functions evaluate to a mean of zero, the expectation of all cross terms strictly vanishes. We are left solely with the expectation of the diagonal terms:
\begin{equation}
    \mathbb{E}\left[|X - Y|^2\right] = \sum_{i=1}^d w_i^2 \mathbb{E}\left[ \big(T_i(N_i) - N_i\big)^2 \right].
\end{equation}
Recognizing that $\mathbb{E}[ (T_i(N_i) - N_i)^2 ]$ is the exact definition of the squared Wasserstein distance between the individual source $Z_i$ and a 
standard normal distribution, we arrive at the final upper bound:
\begin{equation}
    W_2^2(\mathbf{w} \cdot \mathbf{Z}, \mathcal{N}(0,1)) \le \sum_{i=1}^d w_i^2 W_2^2(Z_i, \mathcal{N}(0,1)).
\end{equation}
This convexity adjacent property proves that mixing independent sources mathematically reduces the Wasserstein distance to Gaussianity. 
Maximizing the left-hand side forces the weight vector $\mathbf{w}$ to collapse toward a canonical axis, effectively isolating a single independent source.

\subsection{Comonotonicity and the Proof of Strict Inequality}
While the upper bound demonstrates that mixtures do not exceed the non-Gaussianity of their sources, successfully optimizing ICA requires proving 
that mixtures are strictly less non-Gaussian than the original sources. We must therefore determine the conditions under which the aforementioned 
inequality becomes a strict inequality. 

In a one-dimensional setting, the $W_2$ distance equals the expected squared difference of a specific coupling if and only if the random variables 
are comonotonic. By Brenier's Theorem \cite{brenier1991}, this implies there exists a strictly increasing, deterministic function mapping one variable to the other. 
For our constructed mixture variable $X$ and the Gaussian target $Y$ to be comonotonic with respect to the shared noise vector $\mathbf{N} \in \mathbb{R}^d$, 
their gradients must be perfectly parallel at every point in the probability space. This requires the existence of a scalar scaling function $\lambda(\mathbf{N})$ such that:
\begin{equation}
    \nabla X(\mathbf{N}) = \lambda(\mathbf{N}) \nabla Y(\mathbf{N}).
\end{equation}

Assuming the source distributions possess smooth and strictly positive densities, Caffarelli's regularity theory \cite{caffarelli1992} guarantees that the 
Optimal Transport maps $T_i$ are continuously differentiable. We can therefore compute the gradients of our constructed variables with respect to the noise vector $\mathbf{N}$. 
The gradient of the Gaussian target is simply the weight vector:
\begin{equation}
    Y(\mathbf{N}) = \sum_{i=1}^d w_i N_i \implies \nabla Y = \mathbf{w}.
\end{equation}
The gradient of the source mixture leverages the derivatives of the individual transport maps:
\begin{equation}
    X(\mathbf{N}) = \sum_{i=1}^d w_i T_i(N_i) \implies \nabla X = \begin{bmatrix} w_1 T'_1(N_1) \\ w_2 T'_2(N_2) \\ \vdots \\ w_d T'_d(N_d) \end{bmatrix}.
\end{equation}

Substituting these gradients back into our comonotonicity condition yields the system of equations:
\begin{equation}
    \nabla X = \lambda(\mathbf{N}) \mathbf{w}.
\end{equation}
To determine if this condition can hold for non-Gaussian sources, we differentiate both the Left-Hand Side (LHS) and the Right-Hand Side (RHS) 
with respect to $\mathbf{N}$ to compute their respective Hessian (Jacobian) matrices.

On the left hand side, because $X$ separates into independent terms $w_i T_i(N_i)$, the cross-derivatives $\frac{\partial^2 X}{\partial N_i \partial N_j}$ 
are strictly zero for all $i \neq j$. This yields a purely diagonal matrix containing the second derivatives of the transport maps.

On the right hand side,we differentiate the product $\lambda(\mathbf{N}) \mathbf{w}$. Because the weight vector $\mathbf{w}$ is a constant, 
applying the product rule yields the outer product of $\mathbf{w}$ and the gradient of the scalar function $\nabla \lambda$. This forms a Rank-1 matrix.

Equating the left and right hand sides:
\begin{equation}
    \underbrace{
    \begin{bmatrix} 
    w_1 T''_1(N_1) & 0 & \cdots & 0 \\
    0 & w_2 T''_2(N_2) & \cdots & 0 \\
    \vdots & \vdots & \ddots & \vdots \\
    0 & 0 & \cdots & w_d T''_d(N_d)
    \end{bmatrix}
    }_{\text{Hessian of } X \text{ (Diagonal)}}
    =
    \underbrace{
    \begin{bmatrix}
    w_1 \frac{\partial \lambda}{\partial N_1} & w_1 \frac{\partial \lambda}{\partial N_2} & \cdots & w_1 \frac{\partial \lambda}{\partial N_d} \\
    w_2 \frac{\partial \lambda}{\partial N_1} & w_2 \frac{\partial \lambda}{\partial N_2} & \cdots & w_2 \frac{\partial \lambda}{\partial N_d} \\
    \vdots & \vdots & \ddots & \vdots \\
    w_d \frac{\partial \lambda}{\partial N_1} & w_d \frac{\partial \lambda}{\partial N_2} & \cdots & w_d \frac{\partial \lambda}{\partial N_d}
    \end{bmatrix}
    }_{\text{Outer Product } \mathbf{w} (\nabla \lambda)^\top \text{ (Rank-1)}}
\end{equation}

This matrix equality forces a profound contradiction. Let us examine any off-diagonal entry located at row $i$ and column $j$ (where $i \neq j$). 
The diagonal LHS matrix dictates that this entry must be exactly zero. However, the RHS Rank-1 matrix evaluates this same entry as $w_i \frac{\partial \lambda}{\partial N_j}$. 
This establishes the requirement:
\begin{equation}
    w_i \frac{\partial \lambda}{\partial N_j} = 0 \quad \text{for all } i \neq j.
\end{equation}
Because we are assuming a valid mixture where the component weights are non-zero ($w_i \neq 0$), it must be true that $\frac{\partial \lambda}{\partial N_j} = 0$. 
Since this holds across all dimensions, the gradient of the scaling function is zero ($\nabla \lambda = \mathbf{0}$). 
Consequently, $\lambda(\mathbf{N})$ must be a global scalar constant, $\lambda$.

If $\lambda$ is a constant, then equating the diagonal elements of our matrices implies that $T'_i(N_i) = \lambda$. Integrating this relationship means the 
Optimal Transport maps themselves must be strictly linear functions of the form:
\begin{equation}
    T_i(x) = \lambda x + c.
\end{equation}
We know that applying a purely linear transformation to a Gaussian noise source $N_i$ can only ever produce another Gaussian distribution. 
However, the foundational identifiability assumption of ICA guarantees that our original latent sources $Z_i$ are non-Gaussian. Because the sources are non-Gaussian, 
their Optimal Transport mappings from a Gaussian base must be non-linear (curved). 

Therefore, the linear map requirement derived from the comonotonicity condition fundamentally fails for non-Gaussian sources. 
Because the condition for equality cannot be met, we conclude that for any valid mixture of independent, non-Gaussian variables, the relationship is a strict inequality:
\begin{equation}
    W_2^2(\mathbf{w} \cdot \mathbf{Z}, \mathcal{N}(0,1)) < \sum_{i=1}^d w_i^2 W_2^2(Z_i, \mathcal{N}(0,1)).
\end{equation}

\section{Empirical Landscape: $W_1$ versus $W_2$}
Having theoretically established that the squared Wasserstein distance successfully isolates non-Gaussian sources, we now evaluate its 
practical suitability as an optimization landscape. Specifically, we compare the squared $W_2$ distance (Mean Squared Error ground cost) against the $W_1$ distance 
(Mean Absolute Error ground cost) to observe gradient behavior during rotational search.

\subsection{Student-t Distributions and Degrees of Freedom}
To conduct this empirical comparison, we generated synthetic bivariate mixtures from Student-t distributions. The Student-t distribution is ideal for this analysis 
because its non-Gaussianity can be precisely controlled by adjusting its degrees of freedom ($\nu$). 
A low degrees of freedom parameter (e.g., $\nu=2$) produces heavy tails, representing a strong non-Gaussian signal. 
Conversely, as $\nu$ approaches higher values (e.g., $\nu=10$), the tails become lighter, and the distribution closely approximates a standard normal, 
representing a weak non-Gaussian signal. We systematically evaluated the $W_1$ and $W_2^2$ distances from a standard normal distribution across 
a continuous range of rotation angles $\theta$ relative to the optimal unmixing axis.

\subsection{Robustness to Weak Non-Gaussian Signals}
The empirical results revealed a stark contrast in the optimization landscapes of the two metrics. For strong non-Gaussian signals ($\nu=2$), 
both metrics successfully identified the global optimum, exhibiting distinct peaks aligned with the true source axis. However, the $W_2^2$ landscape 
exhibited significantly higher curvature around the optimum, providing stronger, more directional gradients for convergence. 

The most critical divergence in behavior occurred in the presence of weak signals ($\nu=10$). As the underlying distribution approached Gaussian noise, 
the $W_1$ landscape deteriorated rapidly, becoming highly jagged and non-convex. This roughness introduces severe vulnerability, 
as gradient-based solvers are highly likely to become trapped in spurious local optima. In contrast, the $W_2^2$ metric remained remarkably smooth and unimodal, 
maintaining a clear gradient path toward the true global maximum despite the weak signal strength. This robust empirical behavior, 
combined with the strict mathematical bounds established previously, definitively justifies the computational expense of selecting the 
squared Wasserstein distance as the core contrast function for the proposed W2-ICA framework.

% -----------------------------------------------------------------
% CHAPTER 4: ALGORITHMIC DESIGN AND OPTIMIZATION
% -----------------------------------------------------------------
\chapter{Algorithmic Enhancements}

\section{The Baseline W2-ICA Algorithm}
Having established the theoretical optimality of the squared Wasserstein distance ($W_2^2$) for Independent Component Analysis, 
the challenge transitions to computationally minimizing this cost function. The baseline W2-ICA algorithm proceeds by first centering and whitening 
the observed mixture $\mathbf{X} \in \mathbb{R}^{d \times N}$, yielding the whitened data $\widetilde{\mathbf{X}}$. The algorithm then initializes 
an orthogonal unmixing matrix $\mathbf{W} \in \mathbb{R}^{d \times d}$. In the forward pass, the data is projected onto the candidate components, 
$\mathbf{Y} = \mathbf{W}\widetilde{\mathbf{X}}$. For each row (component) of $\mathbf{Y}$, the empirical cumulative distribution function is 
constructed by sorting the $N$ observations. The $W_2^2$ distance is then computed by matching these sorted empirical quantiles to the corresponding quantiles of a standard normal distribution. 

While theoretically sound, the naive implementation of this gradient descent loop faces severe computational and geometric hurdles. 
The exact Optimal Transport matching requires sorting at every iteration, presenting a significant computational bottleneck. 
Furthermore, standard gradient updates do not inherently respect the strict orthogonality constraint of the mixing matrix, 
and discrete approximations of the target Gaussian introduce gradient noise. To make W2-ICA computationally competitive with existing methods, 
several specific algorithmic and geometric enhancements are required.

\section{Computational Bottlenecks and Solutions}

\subsection{Exact Analytical Gaussian Targets}
A primary source of instability in the baseline algorithm is the discrete approximation of the target Gaussian distribution. 
Traditionally, the target quantiles are approximated by evaluating the inverse cumulative distribution function at evenly spaced intervals, 
such as $q_i = \Phi^{-1}((i - 0.5)/N)$. However, representing a continuous, infinite-tailed Gaussian distribution with a finite set of discrete points 
introduces approximation noise. When the algorithm computes the gradient of the $W_2^2$ distance, this noise translates into a jagged optimization landscape, 
severely hindering convergence as the components approach true independence.

To eliminate this approximation noise, we replace the point-sampled targets with an exact analytical formulation. Using integral calculus, 
we compute the exact expected value of a standard normal variable within each discrete quantile bin. Let the uniform probability bin edges be defined 
as $p_i = \frac{i}{N}$ for $i = 0, \dots, N$, with corresponding Gaussian domain boundaries $z_i = \Phi^{-1}(p_i)$. 
The exact analytical target value $T_i$ for the $i$-th sorted sample is computed as the conditional expectation within that interval:
\begin{equation}
    T_i = N \int_{z_{i-1}}^{z_i} x \phi(x) \, dx
\end{equation}
where $\phi(x)$ is the standard normal probability density function. Evaluating this integral yields the closed-form solution:
\begin{equation}
    T_i = N \big( \phi(z_{i-1}) - \phi(z_i) \big).
\end{equation}
By computing this exact analytical target once during initialization, the algorithm evaluates the $W_2^2$ distance against a mathematically perfect 
discrete representation of a Gaussian, ensuring perfectly smooth gradients near the global optimum.

\subsection{Batched Vectorization for Restarts}
Unlike Principal Component Analysis, which yields a unique global solution via singular value decomposition, 
the ICA optimization landscape is highly non-convex and characterized by numerous local optima. Consequently, 
identifying the true independent components strictly requires multiple random initializations (restarts) to thoroughly explore the hypothesis space. 

Because evaluating the $W_2^2$ objective inherently requires an $\mathcal{O}(N \log N)$ sorting operation for each candidate component, 
executing these random restarts sequentially via standard iterative loops is computationally prohibitive. To resolve this, we utilize a batched vectorization strategy. 
Rather than maintaining a single weight matrix, we aggregate $B$ random restarts into a single high-dimensional tensor $\mathbf{W}_{\text{batch}} \in \mathbb{R}^{B \times d \times d}$. 

This formulation allows the projection $\mathbf{Y} = \mathbf{W}_{\text{batch}} \widetilde{\mathbf{X}}$ to be computed simultaneously for all $B$ trajectories. 
Modern hardware accelerators (GPUs) can then sort the elements of the resulting batch tensor in parallel. 
The pre-computed exact analytical target vector $\mathbf{T}$ is subsequently broadcast-subtracted across the entire batch to compute the distances and gradients in a single step.

\section{Optimization on the Stiefel Manifold}

\subsection{The Riemannian Gradient}
Because the input data is whitened, the unmixing matrix $\mathbf{W}$ must remain strictly orthogonal, satisfying $\mathbf{W}\mathbf{W}^\top = \mathbf{I}$. 
The geometric space of all such orthogonal matrices forms the Stiefel manifold, denoted $\mathcal{S}$. If we compute the standard Euclidean gradient of the squared Wasserstein distance, 
$\mathbf{G} = \nabla_{\mathbf{W}} W_2^2$, this vector points into the flat, unconstrained ambient space $\mathbb{R}^{d \times d}$. 
Updating the matrix using this Euclidean gradient would immediately violate the orthogonality constraint.

To navigate the curved geometry of the Stiefel manifold, we must project the Euclidean gradient onto the tangent space of the manifold at the current point $\mathbf{W}$. 
The resulting Riemannian gradient $\nabla_{\mathcal{S}} \mathbf{W}$ is given by:
\begin{equation}
    \nabla_{\mathcal{S}} \mathbf{W} = \mathbf{G} - \frac{1}{2}\big(\mathbf{G}\mathbf{W}^\top + \mathbf{W}\mathbf{G}^\top \big)\mathbf{W}.
\end{equation}
Geometrically, the term $(\mathbf{G}\mathbf{W}^\top + \mathbf{W}\mathbf{G}^\top)$ precisely measures the symmetric violation of orthogonality induced by the Euclidean gradient. 
By subtracting this specific violation, we ensure that the update step perfectly hugs the curvature of the manifold, preserving the decorrelation of the components.

\subsection{Retraction via Symmetric Decorrelation}
Although the Riemannian gradient ensures the update step is tangential to the Stiefel manifold, 
moving along this flat tangent plane for a finite step size $\eta$ will cause the new matrix $\mathbf{W}_{\text{step}} = \mathbf{W} - \eta \nabla_{\mathcal{S}} \mathbf{W}$ 
to slightly lift off the curved manifold surface. To restore exact orthogonality, the matrix must be mapped back onto the manifold via a mathematical operation known as a retraction.

While the Gram-Schmidt orthogonalization process is a common approach, it is fundamentally asymmetric; 
it permanently fixes the first vector and modifies subsequent vectors to fit, which inherently biases simultaneous optimization. 
Instead, we utilize symmetric decorrelation as our retraction mechanism:
\begin{equation}
    \mathbf{W}_{\text{new}} = (\mathbf{W}_{\text{step}}\mathbf{W}_{\text{step}}^\top)^{-1/2} \mathbf{W}_{\text{step}}.
\end{equation}
This operation computes the overlap covariance of the stepped matrix and applies a uniform, symmetric adjustment to all vectors simultaneously. 
This ensures that no single independent component is arbitrarily prioritized during the retraction process, maintaining the stability of the batched optimization.

\section{The Intractability of Fixed Point Algorithms for W2-ICA}
The industry standard algorithm for ICA, FastICA, achieves highly efficient, super-linear convergence by employing a Newton (second-order) fixed-point step. 
FastICA is mathematically able to utilize this approach because it maximizes proxy contrast functions (such as the \textit{logcosh} function) that are static and easily differentiable. 
A natural question arises: can we formulate a "Fast-Wasserstein-ICA" by deriving an exact Newton step for the $W_2^2$ objective? 

To implement a Newton step, we would require the exact Hessian matrix of the squared Wasserstein distance, $\mathbf{H} = \nabla^2_{\mathbf{w}} W_2^2$. 
The first derivative (the gradient) of the objective with respect to the weight vector $\mathbf{w}$ utilizes the Optimal Transport map $T_{\mathbf{w}}$:
\begin{equation}
    \nabla_{\mathbf{w}} \left( \frac{1}{2} W_2^2 \right) = \mathbb{E}\Big[ \mathbf{X} \big( \mathbf{w}^\top \mathbf{X} - T_{\mathbf{w}}(\mathbf{w}^\top \mathbf{X}) \big) \Big].
\end{equation}
To compute the Hessian, we must differentiate the optimal transport map $T_{\mathbf{w}}(\mathbf{w}^\top \mathbf{X})$ itself with respect to $\mathbf{w}$. 
Applying the total derivative yields two distinct terms:
\begin{equation}
    \nabla_{\mathbf{w}} [T_{\mathbf{w}}(\mathbf{w}^\top \mathbf{X})] = \underbrace{T'_{\mathbf{w}}(\mathbf{w}^\top \mathbf{X}) \mathbf{X}}_{\text{Argument Derivative}} + \underbrace{(\nabla_{\mathbf{w}} T_{\mathbf{w}})(\mathbf{w}^\top \mathbf{X})}_{\text{Map Derivative}}.
\end{equation}
By Caffarelli's regularity theory, the transport map $T_{\mathbf{w}}$ is continuously differentiable, guaranteeing the existence of the Argument Derivative. 
The structural problem lies exclusively within the Map Derivative. 

In the 1D Optimal Transport setting, the transport map to a standard Gaussian relies directly on the empirical cumulative distribution function, $F_{\mathbf{w}}$, 
of the projected data $y = \mathbf{w}^\top \mathbf{X}$. Expanding the Map Derivative yields:
\begin{equation}
    (\nabla_{\mathbf{w}} T_{\mathbf{w}})(y) = \nabla_{\mathbf{w}} \big[ \Phi^{-1}(F_{\mathbf{w}}(y)) \big] = \frac{1}{\phi(\Phi^{-1}(F_{\mathbf{w}}(y)))} \nabla_{\mathbf{w}} F_{\mathbf{w}}(y).
\end{equation}
This introduces a severe mathematical pitfall. The term $\nabla_{\mathbf{w}} F_{\mathbf{w}}(y)$ represents the rate of change of the 
cumulative probability mass as the projection plane $\mathbf{w}$ is rotated. Conceptually and mathematically, the derivative of a 
cumulative distribution function's boundary relies strictly on the underlying Probability Density Function (PDF) evaluated exactly at that boundary. 

This requirement fundamentally defeats the primary advantage of the W2-ICA formulation. 
This advantage of 1D Optimal Transport stems from its reliance on sorting (empirical CDFs), completely circumventing the need for continuous density estimation. 
Requiring the PDF to compute the Hessian forces the algorithm to rely on Kernel Density Estimation (KDE). 
KDE introduces massive statistical noise and bandwidth-tuning dependencies. 

Because the exact Hessian requires continuous density estimation, computing a true Newton-based fixed-point step for exact Wasserstein-ICA is mathematically intractable. 
Consequently, the optimal strategy for the W2-ICA framework is the utilization of Quasi-Newton methods, such as the Limited-memory BFGS (L-BFGS) algorithm adapted for the Stiefel manifold. 
L-BFGS intelligently approximates the inverse Hessian curvature strictly by observing the history of the highly stable, first-order Wasserstein gradients. 
This approach bridges the gap, allowing the algorithm to achieve faster convergence rates while preserving the advantageous, sorting-based methodology of the Optimal Transport formulation.

% -----------------------------------------------------------------
% CHAPTER 5: EXPERIMENTS AND ROBUSTNESS
% -----------------------------------------------------------------
\chapter{Experimental Evaluation and Applications}

\section{The Blind Spots of FastICA}
While FastICA is widely considered the industry standard for independent component recovery, its reliance on fixed proxy contrast functions and a 
Newton-based fixed-point iteration introduces critical mathematical vulnerabilities. In this section, we expose two distinct mathematical blind spots 
that cause FastICA to fail on perfectly valid, non-Gaussian independent components.

\subsection{The Zero-Negentropy Pitfall}
FastICA does not measure true negentropy; instead, it relies on a computationally efficient approximation. For a whitened random variable $x$, 
the negentropy $J(x)$ is approximated using a non-quadratic contrast function $G(\cdot)$, most commonly the \textit{logcosh} function $G(x) = \frac{1}{a} \log \cosh(ax)$. 
The approximation is given by:
\begin{equation}
    J(x) \propto \big( \mathbb{E}[G(x)] - \mathbb{E}[G(\nu)] \big)^2
\end{equation}
where $\nu$ is a standard Gaussian random variable. The algorithm searches for projections that maximize this squared difference. 

The first failure mode occurs if a latent independent source possesses a specific structural shape such that its expected value under the 
contrast function exactly equals that of a Gaussian, i.e., $\mathbb{E}[G(x)] = \mathbb{E}[G(\nu)]$. When this equality holds, 
the approximated negentropy evaluates identically to zero. In this scenario, the objective function completely collapses. 
FastICA becomes entirely blind to the non-Gaussianity of the source, falsely perceiving it as Gaussian noise, and consequently fails to extract it.

\subsection{The Vanishing Curvature Pitfall}
Even if a source successfully bypasses the zero-negentropy pitfall (meaning $J(x) > 0$), the actual \textit{optimization algorithm} used by FastICA can still fail. 
To achieve super-linear convergence, FastICA utilizes a Newton (second-order) fixed-point iteration step. 
However, the mathematical stability of this convergence relies explicitly on a foundational technical condition governing the algorithmic curvature of the optimization landscape. 
As formalized in the FastICA convergence proofs by Hyv{\"a}rinen, Karhunen, and Oja \cite[Chapter 8]{hyvarinen2001}, this requirement is dictated by the following expectation:

The theorem states that for the fixed-point iteration to converge to a true independent component $s_i$, 
the expected value governing the second-order Taylor expansion of the iteration must not vanish. Specifically, defining $g(x) = G'(x)$ and $g'(x) = G''(x)$, the convergence requires:
\begin{equation}
    \mathbb{E}\big[s_i g(s_i) - g'(s_i)\big] \neq 0
\end{equation}
This expectation geometrically represents the algorithmic curvature of the optimization landscape at the optimal point. 
For the standard \textit{logcosh} contrast function (with $a=1$), this equates to $g(x) = \tanh(x)$ and $g'(x) = 1 - \tanh^2(x)$. 

If a specific non-Gaussian source distribution naturally causes this specific expectation to evaluate precisely to zero, 
the algorithmic curvature vanishes. When this condition fails, the denominator of the FastICA update rule collapses, 
destroying the local quadratic (or cubic) convergence guarantees. The fixed-point fails to converge, 
even though the component is strictly independent and non-Gaussian.

\subsection{Constructing the Failure Mode Counterexample}
To empirically demonstrate the severity of the vanishing curvature limitation, we engineered a continuous source distribution specifically designed to trigger the assumption
described in the previous section. 
We analyzed the target function representing the curvature condition: 
\begin{equation}
    h(x) = x \tanh(x) - 1 + \tanh^2(x).
\end{equation}
We observed that $h(x)$ evaluates to negative values near the origin ($|x| < 0.82$) and positive values in the extreme tails. 

We constructed a Trimodal Gaussian Mixture model parameterized by symmetrically spaced peak locations $\{-b, 0, b\}$ and a shared variance $\sigma^2$. 
By solving the integral equation $\int h(x) \text{PDF}(x) dx = 0$ subject to the standard unit variance constraint ($\mathbb{E}[X^2] = 1$), 
we determined the exact probability masses required to perfectly balance the negative central region against the positive tail regions. 
This carefully engineered distribution remains strictly non-Gaussian and perfectly valid under the linear ICA assumptions, 
yet its algorithmic curvature is exactly zero, rendering it mathematically invisible to FastICA's fixed-point solver.

\subsection{Empirical Results: Breaking FastICA}
We evaluated both FastICA and the proposed W2-ICA framework on mixtures of these engineered Trimodal Gaussian sources across 10, 20, and 30 dimensions, 
utilizing a sample size of $N=50,000$. As mathematically predicted by the failure of Assumption A.5, FastICA consistently hit its maximum iteration limits without converging. 
It returned unmixing matrices with catastrophically high Amari errors ($>1.0$), indicating a complete failure to separate the sources.

Conversely, W2-ICA successfully unmixed the 10-dimensional and 20-dimensional signals with high accuracy (Amari error $\approx 0.1$). 
W2-ICA succeeds in this environment because it maps the entire global geometry of the empirical CDF to a Gaussian target via Optimal Transport. 
By utilizing first-order gradient optimization (L-BFGS) on the Stiefel manifold, W2-ICA completely bypasses the local, second-order curvature pitfalls that dismantle Newton-step proxy methods.

\section{Dimensionality Limits and Sample Starvation}
While W2-ICA successfully bypassed the vanishing curvature pitfall, its performance degraded significantly at 30 dimensions with the fixed $50,000$ sample size. 
This performance ceiling is not a structural mathematical failure, but rather an empirical limitation dictated by the Central Limit Theorem.

A linear mixture of 30 independent trimodal sources creates a highly complex joint distribution that appears overwhelmingly Gaussian. 
To mathematically distinguish the microscopic geometric "steps" of this high-dimensional discrete mixture from true continuous Gaussian noise, 
the empirical CDF requires a massive sample size to achieve sufficient resolution. When the sample size is inadequate relative to the dimensionality, 
the system experiences "sample starvation." The subtle non-Gaussian structural signals become overwhelmed by standard empirical sampling noise, 
preventing the Wasserstein gradients from accurately identifying the independent components.

\section{Robust Optimal Transport via Ground Cost Modification}

\subsection{The Outlier Sensitivity of Exact $W_2$}
While the exact squared Wasserstein distance ($W_2^2$) provides superior, smooth gradients for weak signals, it introduces a separate vulnerability. 
Because $W_2$ utilizes an L2 (squared Euclidean) ground cost $c(x,y) = |x - y|^2$, it is notoriously sensitive to outliers, 
such as sensor artifacts or extreme market anomalies. A single massive outlier forces the Optimal Transport plan to expend immense gradient effort to align that single data point, 
massively distorting the overall unmixing matrix and degrading the recovery of the normal data distribution. While $W_2$ is a strict mathematical metric guaranteeing geometric stability, 
it lacks robust statistical properties.

\subsection{The W-Huber Formulation (Logcosh Cost)}
To combine the smooth gradient properties of $W_2$ with the outlier resistance inherent to $W_1$ (which uses an L1 absolute error cost), 
we propose modifying the Optimal Transport ground cost. We replace the strict L2 cost with a statistically robust \textit{logcosh} formulation, creating a "W-Huber" distance metric:
\begin{equation}
    c(x,y) = \log(\cosh(x-y)) \approx |x-y| + \log(1 + e^{-2|x-y|}) - \log(2).
\end{equation}
For small transport distances (representing the core normal data), this cost function behaves quadratically like an L2 penalty, 
preserving the smooth Stiefel gradients necessary for precise convergence. For exceptionally large distances (representing extreme outliers), 
the function transitions to an L1-like absolute penalty, effectively capping the gradient influence of extreme noise. 
Crucially, utilizing the \textit{logcosh} function as a ground cost inside the global Optimal Transport matching process does not trigger the FastICA vanishing curvature pitfall, 
as W-ICA relies on first-order global CDF matching rather than second-order local density approximations.

\section{Application: Price Discovery and Information Shares}

\subsection{Non-Gaussianity in Econometric Market Microstructure}
To demonstrate the empirical utility of the W2-ICA framework, we apply it to the problem of price discovery in financial econometrics. 
In fragmented financial markets where the same asset is traded across multiple venues, 
determining which market contributes most to the incorporation of new information (price discovery) is a non-trivial challenge. 
Standard vector error correction models (VECMs) identify information shares by relying on the uncorrelation of price innovations. 
However, as demonstrated by Zema and Cordoni \citeyear{zema2025}, any identification approach relying solely on uncorrelation is statistically equivalent up to an orthogonal transformation. 

By leveraging the inherent non-Gaussianity of latent market shocks, ICA can be utilized to resolve this orthogonal ambiguity, 
unifying disparate price discovery metrics into a single, unique measurement of market information shares. 

\subsection{Economic Identification Strategy}
While ICA successfully identifies the independent structural shocks driving the market, the raw statistical output remains subject to the inherent permutation and sign ambiguities 
discussed in Chapter 2. To translate these statistical components into meaningful economic drivers, we impose two strict economic identification constraints.

First, to resolve the permutation ambiguity, we enforce a Dominant Diagonal constraint. 
We postulate that the structural shock assigned to Market $i$ must be the specific shock that contributes the maximum variance to that market's price innovations. 
We implement this assignment algorithmically using the Hungarian Algorithm, which uniquely solves the optimal bipartite matching problem between the recovered independent components 
and the individual markets.

Second, to resolve the sign ambiguity, we enforce a Positive Impact constraint based on the economic theory of positive spillovers. 
We assert that a structural innovation representing "good news" must have a positive initial impact on its own market price. 
Any column of the unmixing matrix that violates this economic logic is sign-flipped, ensuring the final structural parameters are economically interpretable. 

\subsection{Empirical Results on Market Data}
[Note: This subsection will be populated with the specific empirical results, tables, and figures generated by the Wasserstein-ICA algorithm applied to the market microstructure data].
% -----------------------------------------------------------------
% CHAPTER 6: CONCLUSION
% -----------------------------------------------------------------
\chapter{Conclusion}
\section{Summary of Contributions}
\section{Limitations}
\section{Future Work}

%%%%%%%%%%%%%%%%%%%%%%%%%%%%%%%%%%%%%%%%%%%%%%%%%%%%%%%%%%%%%%%%%%%
% BACK MATTER
%%%%%%%%%%%%%%%%%%%%%%%%%%%%%%%%%%%%%%%%%%%%%%%%%%%%%%%%%%%%%%%%%%%
\backmatter

\bibliographystyle{apacite}
\bibliography{references} 

\begin{appendices}
    \chapter{Mathematical Proofs and Derivations}
    \section{Derivation of the FastICA Fixed-Point Failure Condition}
    
    \chapter{Additional Algorithms and Code}
\end{appendices}

\end{document}